\section{Study Design}\label{sec:design}

The study measures every rule of Section~\ref{sec:core} against its definitions, asking four questions a developer
has of a slice.

\begin{itemize}
\item \textbf{RQ1} Which dependence rules give sound and dependence-complete slices?
\item \textbf{RQ2} How much smaller than the QDG slice can a slice safely be, and under which inputs?
\item \textbf{RQ3} What does each rule do to fault localization?
\item \textbf{RQ4} What does each rule cost in time as circuits grow?
\end{itemize}

\subsection{Pre-registration}
The protocol was committed before the runs reported here. It extends an earlier run whose results it reuses, so the
outcomes of H1 and H2 were known when it was written, and every later change is a dated deviation. The five
confirmatory tests concern state-based pruning. Every other result is exploratory or descriptive, including the
classical PDG's share of wrong slices and everything added after the first results, the classical PDG and
\texttt{LightCone} baselines, the dependence set and its nesting with the rules, the symbolic-angle, flow-sensitive and
read-only variants, and the gate check. Two descriptive tasks for language models were pre-registered, and the most
expensive model was dropped for cost before any scored call.

\subsection{Subjects}
The rules are compared on circuits, where exact simulation gives the ground truth. Corpus~A collects
\numSubjACount{} benchmark circuits from MQT Bench~\cite{quetschlich2023mqt}, QASMBench~\cite{li2023qasmbench} and
VeriQBench~\cite{chen2022veriqbench}, with \numSubjAMinQubits{} to \numSubjAMaxQubits{} qubits and a median of
\numSubjAMedianOps{} operations, at most \numSubjAMaxOps{}. Every qubit of every circuit is a criterion, which gives
\numNPairs{} circuit--qubit pairs, \numNDynamicPairs{} of them in dynamic circuits with mid-circuit measurement, reset
or conditioned gates. The \numSubjGenCount{} generated circuits come in families of known structure that grow in depth
and in width, with \numSubjGenMinQubits{} to \numSubjGenMaxQubits{} qubits and up to \numSubjGenMaxOps{} operations,
and \numSubjSweepCount{} width sweeps extend two of the families to \numSubjSweepMaxQubits{} qubits. Corpus~L has
\numSubjLCount{} wider circuits from the same suites, with \numSubjLMinQubits{} to \numSubjLMaxQubits{} qubits and up
to \numSubjLMaxOps{} operations, on which only some rules finish. Faults are seeded into corpus~A, and real faults
come from Bugs4Q~\cite{zhao2023bugs4q}. A unit is a fault and a qubit whose measured distribution it changes.

\subsection{Ground truth}
An independent oracle simulates each circuit exactly, branching over measurement outcomes, and agrees with Qiskit's
simulation to within \numOracleErr{}. A slice preserves the criterion when the squared Uhlmann fidelity of the
criterion's reduced states under the slice and the whole circuit is at least $1 - 10^{-6}$. The
dependences of Definition~\ref{def:dep} are measured by inserting \numInfluencePerturbations{} random unitaries after
every operation and comparing the criterion's state at \numInfluenceInputs{} inputs, $\ket{0\dots0}$ and random
product states. Being sampled, the measured dependence set can only miss true dependences. A catalog of \numCatEntries{} hand-derived cases checks the ground truth and every
slicer before the runs.

\subsection{Slicers}
Each rule is implemented on circuits from its read and write sets or its pruning step. Qiskit's \texttt{LightCone}
pass applies to \numLcqApplicable{} of the \numNPairs{} pairs. It rejects mid-circuit measurements on
\numLcqNaMeasurements{} pairs, and on \numLcqNaReorder{} its output cannot be matched back to the original.
Exploratory variants make commutation pruning hold for all angles and give state-based pruning a joint criterion.

\subsection{Robustness}
A slice computed at $\ket{0\dots0}$ is rechecked, with its operations unchanged, under \numRobDraws{} random
computational-basis inputs, \numRobDraws{} random product inputs and, for circuits with rotation angles,
\numRobDraws{} random rebindings of every angle. A pair is robust if the slice preserves the criterion in every draw. The gate check takes each of the
\numGateSignalN{} two-qubit gates of Qiskit's standard library, parameterized gates at a fixed angle, and searches
seeded random product inputs, for each direction, for two inputs of one qubit that leave the other in different
reduced states, evidence for these gates and not a proof for all.

\subsection{Faults and fault localization}
Faults are seeded with \numMutOperators{} operators, \numMutPublished{} from Muskit~\cite{mendiluze2021muskit},
QMutPy~\cite{fortunato2022qmutpy} and García de la Barrera et al.~\cite{garcia2022testing}, and an angle perturbation
of our own. Of the \numMutants{} mutants, \numMutantsExcluded{} with an empty fault site are excluded, and
\numMutantsKilled{} change the measured distribution of at least one qubit. Each such qubit and its mutant form a
unit, \numRqThreeUnits{} in all, for which we record each slice's size and whether it contains the fault site. Of the
\numBugsFourQTriaged{} entries of Bugs4Q, \numBugsFourQIncluded{} are circuit-level faults that port mechanically to
the current Qiskit release, and \numBugsFourQKilled{} of them change a qubit's measured distribution, giving
\numLinesUnits{} units. Running each port
records the line that appends each operation, mapped to the original program, so these slices are also reported as
lines of the programs their developers wrote.

Mutation-based fault localization ranks the operations of the faulty circuit by how their mutants change its
observable behavior, with the Ochiai formula, as Ishimoto et al.\ do for quantum programs~\cite{ishimoto2025mbfl}, on a
sample of \numMbflUnits{} units stratified by fault operator and including every Bugs4Q unit. We rank over the whole
circuit and within each slice and report the EXAM score, the share of ranked operations examined before the fault,
lower being better.

In the language-model tasks, a model lists the operations that can affect a qubit for \numTOneItems{} items, and
names the likely faulty operation of \numTTwoItems{} units given the whole circuit or a slice. The models are Claude
Sonnet~5.5 and Haiku~5.5 and Qwen2.5-Coder with 32 and 1.5 billion parameters.

\subsection{Cost}
RQ4 times the rules on the generated family that grows in depth at eight qubits, on the width sweeps and on
corpus~L, each circuit in a fresh process limited to \numCostWallLimit{}~s and 4~GiB.

\subsection{Statistics}
Five confirmatory hypotheses concern state-based pruning. It is wrong on fewer pairs than def-use (H1), its slices
are smaller than QDG slices (H2), at least 95\% of pairs keep a correct slice under random product inputs (H3), it
loses at most five points of fault containment against the QDG slice (H4), and its slices are smaller where both
contain the fault (H5). It starts from the QDG slice and checks every removal at $\ket{0\dots0}$, so H1 and the
directions of H2 and H5 hold by construction, and H3 and H4 are the real tests. We use exact McNemar tests on paired outcomes and Wilcoxon signed-rank tests on paired
sizes, with Holm's correction across the five. Effect sizes come with bootstrap 95\% intervals and Wilson intervals
for proportions, and exploratory intervals resample whole algorithm families. Everything else is descriptive.
